\section{Introduction}
In clinical Magnetic Resonance Imaging (MRI), reducing scan time is crucial for patient comfort and minimizing motion artifacts \cite{lustig2007sparse}. However, accelerating MRI scans by undersampling in k-space inevitably leads to severe aliasing artifacts in the reconstructed images. To recover high-quality, aliasing-free images while maintaining high acceleration factors, deep learning approaches have shown tremendous potential and superiority over traditional compressed sensing methods \cite{schlemper2017deep, wang2016accelerating}. This project aims to utilize the BraTS (Brain Tumor Segmentation) dataset to develop a robust, deep learning-based MRI reconstruction algorithm.

In the initial phase (Task 2), we implemented a baseline reconstruction network employing a standard U-Net architecture \cite{ronneberger2015u}. This model was trained using an L2 loss (Mean Squared Error) objective to minimize the pixel-wise discrepancy between the reconstructed image and the fully sampled ground truth. However, an in-depth analysis of the baseline results revealed that relying solely on L2 loss constraints inherently causes the network to generate over-smoothed images, resulting in the loss of critical high-frequency edges and fine texture details during reconstruction \cite{johnson2016perceptual}.

More importantly, we observed a significant limitation during evaluation: reconstructed images with a higher Peak Signal-to-Noise Ratio (PSNR) do not necessarily exhibit better image quality from a human visual perspective. To thoroughly investigate and validate this phenomenon, drawing inspiration from studies on image sharpness, we designed and conducted a ``shifting pixels'' experiment. The experimental results compellingly demonstrated that when an image undergoes a minor spatial shift, the PSNR drops drastically, yet the subjective image quality perceived by the human eye remains completely unchanged. This underscores the severe flaws of traditional pixel-wise metrics in reflecting true visual fidelity \cite{ledig2017photo}.

Based on these observations, we made the following core contributions in the advanced reconstruction phase of this project:
\begin{itemize}
    \item \textbf{Integration of Wavelet Loss:} To address the over-smoothing issue in the U-Net baseline, we introduced a wavelet loss constraint in the frequency domain. This structural penalty enables the network to better focus on and recover the high-frequency details of the image \cite{kang2018deep}.
    \item \textbf{Introduction of DISTS Metric and Perceptual Loss:} To tackle the misalignment between PSNR and human visual perception, we incorporated the Deep Image Structure and Texture Similarity (DISTS) metric into our evaluation framework, which explicitly focuses on texture and structural fidelity \cite{ding2020image}. Concurrently, we integrated a perceptual loss into our training objectives \cite{johnson2016perceptual}. This guides the network to approximate the ground truth at the feature level rather than merely the pixel level, significantly enhancing subjective visual quality.
    \item \textbf{Multi-modal Fusion and Unrolled Reconstruction:} To achieve clinical-grade quality, we developed a sophisticated multi-modal fusion architecture in Task 3. This design takes two modalities as input: an undersampled T2 image and a fully sampled T1 image, leveraging the anatomical structural prior from T1 to help recover T2 details \cite{zhou2020dudornet}. Furthermore, we implemented an unrolled reconstruction network by cascading the denoising network with Data Consistency (DC) layers to strictly ensure k-space fidelity \cite{schlemper2017deep}.
\end{itemize}

Ultimately, we seamlessly integrated the wavelet loss and perceptual loss into the training of the multi-modal unrolled network. Experimental results prove that our complete pipeline not only achieves excellent performance in quantitative metrics but also makes significant breakthroughs in preserving tumor details and eliminating aliasing artifacts visually.




